\documentclass[conference]{IEEEtran}

% *** PACKAGES ***
\usepackage[utf8]{inputenc}
\usepackage[T1]{fontenc}
\usepackage{cite}
\usepackage{graphicx}
\usepackage{amsmath,amssymb}
\usepackage{booktabs}
\usepackage{url}
\usepackage{xcolor}

\begin{document}

% *** TITLE & AUTHORS ***
\title{A Survey of Fine-Tuning Frameworks and Reinforcement Learning Tooling for Speech-to-Text}

\author{
    \IEEEauthorblockN{
    Gautam Agarwal\IEEEauthorrefmark{1},
    Shivangi Kumar\IEEEauthorrefmark{1},
    Chen Wang\IEEEauthorrefmark{2}}
    \IEEEauthorblockA{\IEEEauthorrefmark{1}Department of Computer Science, Columbia University, New York, USA\\
    Email: \{first.last\}@example.edu}
    \IEEEauthorblockA{\IEEEauthorrefmark{2}IBM Research Center, New York, USA\\
    Email: second@example.com}
}

\maketitle

% *** ABSTRACT & INDEX TERMS ***
\begin{abstract}
% TODO: 150--250 word summary of:
% - Motivation: domain shift, static STT deployments, lack of self-adaptation
% - Scope: survey of STT fine-tuning toolkits (NeMo, ESPNet, Align-SLM / SpeechLM)
%          and RL/RLHF frameworks used for model alignment
% - Contributions: taxonomy, framework analyses, experiments, and gap analysis
% - Outcome: design requirements for self-adaptive STT, to be addressed in follow-up work
\end{abstract}

\begin{IEEEkeywords}
Speech-to-Text, Reinforcement Learning, RLHF, Continuous Learning, Fine-Tuning, NeMo, ESPNet, Align-SLM
\end{IEEEkeywords}

% *** MAIN BODY ***

\section{Introduction}
% A. Motivation and Problem Context
% - Degradation of STT under domain shift, noise, accents
% - Limitations of static training + periodic manual retraining
% - Lack of closed-loop, self-improving STT systems

% B. Scope and Survey Objectives
% - Survey STT fine-tuning frameworks: NeMo, ESPNet, Align-SLM / SpeechLM toolkits
% - High-level scan of RL/RLHF frameworks relevant to STT (VERL, TRL, DeepSpeed-Chat, etc.)
% - Focus on identifying design gaps rather than proposing a full solution

% C. Contributions and Paper Organization
% - Taxonomy of STT fine-tuning frameworks
% - Mapping between STT toolchains and RL/RLHF components
% - Experimental comparison on shared datasets and metrics
% - Gap analysis and research opportunities, motivating follow-up work

\section{Background and Problem Setting}
% A. Modern STT Architectures and Supervised Fine-Tuning
% - Brief overview of encoder--decoder, CTC, RNN-T, Conformer-based models
% - Standard supervised fine-tuning loop and metrics (WER, CER, etc.)

% B. Reinforcement Learning and RLHF for Sequence Models
% - RL basics for sequence modeling (policies, rewards, trajectories)
% - RLHF pipeline: SFT, reward modeling, PPO/GRPO/DPO, RLAIF

% C. Continuous and Lifelong Learning in STT
% - Catastrophic forgetting, backward transfer, replay buffers
% - Scheduling problem: when to fine-tune under cost constraints

\section{Survey Methodology and Evaluation Protocol}
% A. Framework Selection Criteria
% - Inclusion criteria for STT toolkits and RL/RLHF frameworks
% - Versioning and documentation sources

% B. Experimental Protocol for STT Frameworks
% - Datasets: list and characteristics (clean, noisy, domain-specific)
% - Metrics: WER, CER, DER, Verb Error Rate, Domain Error Rate, latency, cost

% C. Reproducibility Setup
% - Docker images, GCP or cluster configuration (high level)
% - Random seeds, hyperparameter ranges

\section{Taxonomy of STT Fine-Tuning Frameworks}
% A. Taxonomy Dimensions
% - Model support (CTC, RNN-T, encoder--decoder, streaming/offline)
% - Adaptation strategies (full FT, PEFT, multi-stage training)
% - Data/recipe ecosystems and configuration styles

% B. RL Alignment and Control Hooks as a Taxonomy Dimension
% - Where reward signals can be injected (loss abstraction, scorers)
% - Modularity of training loops; logging, replay, online/incremental support

% (Optional) C. Summary Table
% - Table: frameworks vs. taxonomy dimensions
% \begin{table}[t]
% \centering
% \caption{Taxonomy of STT Fine-Tuning Frameworks}
% \begin{tabular}{lccc}
% \toprule
% Framework & Model Types & PEFT & RL Hooks \\
% \midrule
% NeMo     & ... & ... & ... \\
% ESPNet   & ... & ... & ... \\
% Align-SLM / SpeechLM & ... & ... & ... \\
% \bottomrule
% \end{tabular}
% \label{tab:taxonomy}
% \end{table}

\section{NVIDIA NeMo for STT Fine-Tuning}
% A. Framework Overview
% - ASR pipeline: data loaders, models, loss functions, config-driven training

% B. Fine-Tuning Strategies
% - Full fine-tuning vs LoRA/adapters/BitFit and their parameter efficiency
% - Domain/speaker adaptation examples

% C. RL/RLHF Compatibility Analysis
% - Existing RLHF support (NeMo-Aligner for LLMs)
% - Feasibility of integrating reward-based objectives for STT
% - Limitations for online and streaming RL

\section{ESPNet and ESPnet-SpeechLM}
% A. ESPNet Overview
% - ASR and broader speech toolkit; recipe structure, hybrid CTC/attention models

% B. Fine-Tuning and Adaptation Mechanisms
% - Domain/speaker adaptation recipes; configuration hooks

% C. RLHF and Preference-Based Learning
% - ESPnet-SpeechLM and RLHF capabilities
% - Applicability to STT and gaps (WER, streaming, deployment constraints)

\section{Align-SLM as a Speech RL Alignment Case Study}
% A. Align-SLM Overview
% - Textless spoken language modeling
% - RLAIF / DPO training with AI-generated preferences

% B. Lessons for STT
% - What transfers: preference-based alignment, semantic metrics, AI feedback
% - What does not: generative continuation vs. transcription, CTC alignment

\section{Survey of RL/RLHF Frameworks Relevant to STT}
% A. Taxonomy of RL/RLHF Frameworks
% - Groups: LM-focused RLHF toolkits vs general RL libraries

% B. High-Level Summaries
% - VERL, AREAL, MILES, SMILE
% - OpenRLHF, TRL, DeepSpeed-Chat, TRLX, RL4LMs
% - Ray RLlib and related generic RL libraries

% C. STT Integration Gap (Teaser)
% - Lack of native audio/CTC environments, replay tailored to STT,
%   and tight coupling to NeMo/ESPNet

\section{Cross-Framework Mapping: STT vs. RL/RLHF Components}
% A. Conceptual Mapping
% - STT components: data loaders, decoders, trainers, metrics
% - RL/RLHF components: environment, policy, reward, replay buffer, optimizer

% B. Integration Pain Points
% - Training loop ownership, observation/action mismatch
% - Performance and deployment constraints for STT vs text LMs

% (Optional) C. Mapping Table
% - Framework-by-framework mapping of available/missing pieces

\section{Experimental Comparison of STT Fine-Tuning Frameworks}
% A. Experimental Design
% - Baseline fine-tuning for NeMo, ESPNet (and others if applicable)
% - PEFT variants, multi-domain and drift scenarios

% B. Metrics and Evaluation
% - Accuracy, robustness under domain shift, continual-learning metrics
% - Efficiency: latency, GPU-hours, cost per 1\% WER improvement

% C. Results and Discussion
% - Summary of key trends; framework strengths/weaknesses from experiments
% - Observations that feed directly into the gap analysis

\section{Gap Analysis and Research Opportunities}
% A. Gaps in STT Fine-Tuning Frameworks
% - Missing online/on-policy adaptation interfaces
% - Limited continual-learning and replay-buffer support
% - Lack of standardized reward interfaces over STT metrics

% B. Gaps in RL/RLHF Frameworks for STT
% - Text-only assumptions, no CTC/ASR awareness
% - No turnkey integration with NeMo/ESPNet

% C. Design Requirements for Self-Adaptive STT
% - Desired properties of a unified, self-learning STT framework
% - Keep description abstract to avoid overlapping with the follow-up paper

\section{Open Challenges and Future Directions}
% A. Technical Challenges
% - Sample efficiency and stability of RL/RLHF for STT
% - Robustness to noisy or adversarial feedback
% - Multi-objective optimization (accuracy, latency, cost, safety)

% B. Future Research Directions
% - Preference-based alignment for STT
% - Agentic STT systems with autonomous data curation and scheduling
% - Integration into production MLOps and monitoring pipelines

\section{Conclusion}
% - Concise recap of survey findings
% - Reiterate that existing tools fall short for full self-adaptive STT
% - Explicitly state that this motivates dedicated frameworks,
%   to be addressed in subsequent work

% *** ACKNOWLEDGMENT (OPTIONAL) ***
\section*{Acknowledgment}
% TODO: Acknowledge funding, collaborators, data providers, etc.

% *** REFERENCES ***
\bibliographystyle{IEEEtran}
\bibliography{references} % Create references.bib and manage with BibTeX

\end{document}